\documentclass[11pt,letterpaper]{article}
\usepackage{fontspec}
\usepackage{tocloft}
\usepackage{multicol}
\usepackage{fancyhdr}
\usepackage{nicefrac}
\usepackage[left=1.3in,right=1.3in,top=1in,bottom=1in]{geometry}
\setmainfont[Scale=1.3,AutoFakeBold=1.5,AutoFakeSlant=0.3]{Doves Type}

\widowpenalty=10000
\clubpenalty=10000
\interlinepenalty=500

\title{Escoffier Brown Stock}
\author{}
\date{}

\pagestyle{fancy}
\fancyhf{}
\fancyhead[C]{\textit{Escoffier Brown Stock}}
\fancyfoot[C]{\thepage}
\renewcommand{\headrulewidth}{0pt}
\renewcommand{\footrulewidth}{0pt}

\begin{document}

\maketitle
\thispagestyle{empty}

\begin{quote}
\textit{Classic French brown stock in two stages: beef bones are roasted until dark brown, then simmered into a bone broth for 5--6~hours; the broth is used to deglaze seared stewing meat in repeated reductions until syrupy, then the remaining broth and meat simmer together for about 2~hours. The result is an intense, gelatinous stock for building sauces and braises. Yields about 8~cups.}
\end{quote}

\section*{Ingredients}

\setlength{\columnsep}{20pt}
\begin{multicols}{2}
\noindent
    Beef bones, cut into chunks \dotfill 5--5\nicefrac{1}{2}~lb. \\
    Carrots \dotfill 5~oz. (about 1~cup diced) \\
    Onion \dotfill 5~oz. (about 1~cup diced) \\
    Fresh parsley stalks \dotfill 5 \\
    Fresh thyme sprigs \dotfill 2 \\
    Fresh bay leaf \dotfill 1 \\
    Garlic clove \dotfill 1 \\
    Cold water \dotfill 21~cups \\
    \columnbreak
    Beef or veal stewing meat (knuckle, shank, or chuck), boned and cut into large dice \dotfill 2~lb. \\
    Ham knuckle (small) \dotfill 1 \\
    Fresh pork rind, blanched \dotfill 5~oz. \\
    Neutral oil \dotfill 1--2~Tbsp. \\
\end{multicols}

\section*{Directions}

\noindent
Have butcher cut \textbf{beef bones} into chunks (or crack at home) ---
Make a bouquet garni: tie \textbf{parsley}, \textbf{thyme}, and \textbf{bay leaf} together with kitchen string (\textit{Small Bowl~\#1}) ---
Roughly dice \textbf{carrots} and \textbf{onion}; combine in \textit{Small Bowl~\#2} (aromatics) ---
Cut \textbf{stewing meat} into large dice and set aside in \textit{Large Bowl~\#1}

\begin{enumerate}
    \item Preheat oven to \textit{450--475°F}. Spread \textbf{beef bones} in a single layer in a large roasting pan. Roast until deeply browned, \textit{45--60~minutes}, turning once or twice so they color evenly. Bones are done when they are dark brown all over and smell roasted; avoid charring. If bones release a lot of fat, you can pour off excess partway through.

    \item Add \textbf{carrots} and \textbf{onion} (\textit{Small Bowl~\#2}) to the roasting pan with the bones and stir to coat in fat. Return to the oven and cook for \textit{5~minutes} until vegetables are lightly seared and fragrant.

    \item Remove the pan from the oven. Set a sieve over a bowl and transfer \textbf{roasted bones} and \textbf{vegetables} into the sieve; let drip for \textit{a few minutes} to drain excess fat. Discard or reserve the fat; transfer bones and vegetables to a large stockpot (at least 12~qt.). Add the \textbf{bouquet garni} (\textit{Small Bowl~\#1}) and \textbf{garlic clove}. Pour in 17~cups cold \textbf{water}---the water should sit just a few inches above the bones. Bring to a boil over high heat.

    \item Once boiling, skim the scum that rises to the surface for \textit{2--3~minutes}. Reduce heat to \textit{low} so the liquid maintains a bare simmer---small bubbles occasionally breaking the surface, not a rolling boil. Simmer for \textit{5--6~hours}, skimming any foam or impurities from the surface as needed and topping up with additional \textbf{water} (about 4~cups total) so the liquid level stays roughly constant. The bone broth is ready when it is deeply colored and rich-tasting; use it when you are ready to start Phase~2.

    \item \textbf{Phase~2:} In a second large stockpot or very large saucepan (at least 8~qt.), heat 1--2~Tbsp. \textbf{oil} over \textit{medium-high heat}. Add \textbf{diced stewing meat} (\textit{Large Bowl~\#1}) in a single layer (work in batches if needed to avoid crowding). Sear until lightly browned on multiple sides, \textit{6--10~minutes} per batch. Meat is done when surfaces are browned and no raw gray remains on the outside; it need not be cooked through.

    \item Pour about 1~cup of the \textbf{bone broth} from the first stockpot over the seared meat. Bring to a simmer over \textit{medium heat} and cook until the liquid has almost completely evaporated and reduced to a syrupy glaze coating the bottom of the pot, \textit{10--15~minutes}. Repeat this deglazing and reduction with another 1~cup broth, then a third time. After the third reduction, the bottom should be coated in a dark, sticky glaze.

    \item Add the remaining \textbf{bone broth} from the first stockpot to the pot with the meat and glaze. Add the \textbf{ham knuckle} and \textbf{blanched pork rind}. Bring to a boil, then skim any impurities from the surface. Reduce heat to \textit{low} and simmer gently for \textit{about 2~hours}, until the stock is deeply flavored and the meat has given up its flavor. The stock should remain at a bare simmer; adjust heat as needed.

    \item Turn off the heat and let the stock rest for \textit{at least 30~minutes}. Set a fine-mesh sieve (and optionally a second pass with cheesecloth) over a clean large container. Ladle or pour the stock through the sieve, leaving the solids behind. Do not press the solids; let the stock drain naturally for a clearer result.

    \item Let the strained stock cool to room temperature, then refrigerate uncovered or loosely covered until cold, \textit{overnight or at least 4~hours}. A firm cap of fat will form on the surface. Scoop off the fat and reserve for cooking or discard. Use the stock within \textit{4--5~days} refrigerated or freeze in portions for \textit{up to 3~months}.
\end{enumerate}

\newpage

\section*{Equipment Required}
\begin{itemize}
    \item Large roasting pan (enough to hold bones in one layer)
    \item Large stockpot, at least 12~qt. (for bone broth)
    \item Second large stockpot or saucepan, at least 8~qt. (for Phase~2)
    \item Fine-mesh sieve
    \item Cheesecloth (optional, for second straining)
    \item Large heatproof container or bowl for strained stock
    \item Ladle
    \item Kitchen string (for bouquet garni)
    \item Cutting board and knife
\end{itemize}

\section*{Hints and Notes}

\subsection*{Yield}
\begin{itemize}
    \item Yields about 8~cups brown stock
    \item Use as base for French sauces (espagnole, demi-glace), braises, and soups
\end{itemize}

\subsection*{Mise en Place}
\begin{itemize}
    \item Small Bowl \#1 --- bouquet garni: \textbf{parsley}, \textbf{thyme}, \textbf{bay leaf} tied with string
    \item Small Bowl \#2 --- aromatics: roughly diced \textbf{carrots} and \textbf{onion} (about 1~cup each)
    \item Large Bowl \#1 --- \textbf{stewing meat} cut into large dice (about 2~lb.)
    \item Have \textbf{ham knuckle} and \textbf{blanched pork rind} ready before Phase~2
    \item Bone broth is used directly from the first stockpot; no separate bowl needed
\end{itemize}

\subsection*{Ingredient Tips}
\begin{itemize}
    \item Ask the butcher to cut \textbf{beef bones} into 2--3-inch chunks (marrow bones, knuckle, neck, or mixed); cracking exposes marrow and improves extraction
    \item \textbf{Stewing meat}: use gelatin-rich cuts such as beef shank, veal knuckle, or chuck; avoid lean cuts
    \item \textbf{Ham knuckle} (small ham hock) adds depth and a touch of smoke; find at butchers or well-stocked grocers
    \item \textbf{Blanched pork rind}: buy from a butcher or Asian market; blanch in boiling water \textit{2--3~minutes}, drain, and scrape any remaining hair before use. It contributes body and gelatin.
\end{itemize}

\subsection*{Preparation Tips}
\begin{itemize}
    \item Do not fry bones and vegetables together in fat before adding water---Escoffier's method relies on roasting bones only, then diffusion for color; pre-frying risks over-coloring and bitter flavor
    \item Maintain a bare simmer, not a boil, for both phases to keep the stock clear and avoid emulsifying fat and protein
    \item Skim consistently during the first phase; less scum forms in Phase~2 but skim when it appears
    \item The three deglazes in Phase~2 build a fond that dissolves into the stock for intensity; reduce each addition until truly syrupy before adding the next
    \item Resting \textit{30~minutes} before straining lets solids settle and yields a clearer stock
\end{itemize}

\subsection*{Make Ahead \& Storage}
\begin{itemize}
    \item Stock keeps refrigerated \textit{4--5~days}; freeze in 1--2~cup portions for \textit{up to 3~months}
    \item Cool to room temperature before refrigerating; do not seal tightly while still warm
    \item Remove the fat cap after chilling for easier scooping; reserved fat can be used for roasting or sautéing
    \item Thaw frozen stock in the refrigerator or in a cool water bath; reheat gently
\end{itemize}

\subsection*{Serving Suggestions}
\begin{itemize}
    \item Base for brown sauces (espagnole, demi-glace), pan sauces, and gravies
    \item Liquid for braised beef, pot roast, and daube
    \item Soup base for French onion soup, consommé, or vegetable-beef soup
    \item Reduce further for glazes or au jus
\end{itemize}

\end{document}
